Testing began with a set of candidate basic questions, with the goal of getting the system working before tackling more complex cases. The results were immediately impressive: the required operators handled simple filter-aggregate-sort pipelines swiftly and cleanly. As more advanced questions were introduced and table support operators -- snapshot, restore, and join -- were added to handle dataset-wide aggregation, the system remained surprisingly versatile. The core challenge with advanced questions was that aggregation inherently destroys information, collapsing columns that downstream steps expected to still be present. Table snapshots solved this by allowing the planner to checkpoint the full table before an aggregation and rejoin it afterward, restoring the columns needed for subsequent steps. Beyond this pattern, the required operator set proved expressive enough to cover all fifteen questions without needing to force or workaround any of them.

The most fundamental limitation of the current design is its unsuitability for inference -- drawing conclusions that go beyond what is directly computable from the data. This is not entirely a weakness; it seems improper for an analytical agent to produce answers not clearly grounded in the data. But it does mean the system cannot, for example, identify trends, generate predictions, or reason about causality. Supporting anything like that would require a substantially more complex architecture: replanning, sub-task decomposition, and tighter feedback between the planner and executor. One concrete extension worth exploring is a flexible escape-hatch operator that delegates to short LLM-generated code when no existing operator fits. This would reintroduce some of the expressiveness of the HW1 agent in a controlled way, and it would be informative to observe which computations the planner chose to offload.

Several architectural improvements shaped the current system and point toward future directions. The retry mechanism is functional but shallow -- on failure, the system simply replans from scratch without passing any error context back to the planner. A meaningful improvement would be structured error propagation, so the planner understands what went wrong and can correct the specific step rather than regenerating the entire plan. One change that proved genuinely impactful was having the planner produce a description of the data it expected the executor to return; this context is passed to the response node, which produces noticeably more grounded and accurate final answers. In a deployed interface, surfacing the plan and its reasoning to the user before execution would further improve transparency -- letting users catch planning errors early rather than only seeing a final answer that may be the result of a subtly incorrect plan.
